\documentclass[a4paper,11pt]{article}
        \usepackage[top=15mm,bottom=18mm,left=16mm,right=16mm]{geometry}
        \usepackage{fontspec}
        \usepackage{xeCJK}
        \setmainfont{Times New Roman}
        \setCJKmainfont{Yu Gothic}
        \setmonofont{Consolas}
        \setCJKmonofont{Yu Gothic}
        \usepackage{booktabs}
        \usepackage{longtable}
        \usepackage{tabularx}
        \usepackage{array}
        \usepackage{enumitem}
        \usepackage{hyperref}
        \usepackage{listings}
        \usepackage{xcolor}
        \usepackage{amsmath}
        \usepackage{amssymb}
        \hypersetup{
          colorlinks=true,
          linkcolor=blue,
          urlcolor=blue,
          pdftitle={上位MPC 目的関数},
          pdfauthor={Codex}
        }
        \lstset{
          basicstyle=\ttfamily\small,
          breaklines=true,
          columns=fullflexible,
          keepspaces=true,
          frame=single,
          framerule=0.2pt,
          rulecolor=\color{black},
          showstringspaces=false
        }
        \setlength{\parskip}{0.42em}
        \setlength{\parindent}{1em}
        \setlength{\emergencystretch}{3em}
        \renewcommand{\arraystretch}{1.15}
        \title{16. 上位MPC 目的関数}
        \author{solar\_ws0129-main}
        \date{2026-07-29}
        \begin{document}
        \maketitle
        \tableofcontents
        \clearpage

        \section{対象}
        \begin{tabularx}{\linewidth}{>{\raggedright\arraybackslash}p{0.18\linewidth} X}
        \toprule
        項目 & 内容 \\
        \midrule
        ファイル & \path|mpc_solarcar/upper_cost.py| \\
        ソースSHA-256 & \path|b69065249298ad089d085bcb06d9277c10dae6b61ec0d334d69a13ffef56a3bb| \\
        種別 & Python \\
        区分 & planner helper \\
        役割 & 速度、SoC、温度、電流、進捗、day-end、terminal 条件を penalty として定義する。 \\
        起動文脈 & upper planner の良し悪しを数式化する場所。 \\
        \bottomrule
        \end{tabularx}

        \section{このファイルを読む理由}
        速度、SoC、温度、電流、進捗、day-end、terminal 条件を penalty として定義する。

        \begin{itemize}[leftmargin=1.6em]
        \item upper\_stage\_cost と upper\_terminal\_cost が中心。
\item load\_upper\_cost\_config が profile の重み束を解く。
        \end{itemize}

        \section{呼び出し関係}
        \subsection{どこから呼ばれるか}
        \begin{itemize}[leftmargin=1.6em]
        \item \path|mpc_node.py|
\item \path|scripts/solar_sim.py|
\item \path|tune_upper_planner_weights.py|
        \end{itemize}

        \subsection{次に読むと理解が進むファイル}
        \begin{itemize}[leftmargin=1.6em]
        \item \path|scripts/tune_upper_planner_weights.py|
        \end{itemize}

        \subsection{同一リポジトリ内の主要依存}
        \begin{itemize}[leftmargin=1.6em]
        \item 同一リポジトリ内の明示 import は少ない。
        \end{itemize}

        \section{ソース冒頭の把握}
        モジュール docstring または先頭意図の要約:

        モジュール docstring は明示されていない。

        \subsection{代表コード断片}
        \begin{lstlisting}[language=Python]
@dataclass
class UpperCostConfig:
    objective_mode: str = "weighted"
    w_wait: float = 1.0
    w_travel_time: float = 1.0
    w_terminal_soc_min: float = 30.0
    w_day_end_soc_min: float = 1.0e5
    w_soc_day_max: float = 1.0e4
    w_soc_day_track: float = 0.0
    w_speed_smooth: float = 30.0
    w_dv_limit: float = 2.0
    w_speed_limit: float = 50.0
    w_drive_window: float = 1.0e5
    w_current_sq: float = 0.01
    w_pack_energy: float = 0.0
    w_joule_loss: float = 0.0
    w_aero_energy: float = 0.0
    w_mech_energy: float = 0.0
    w_speed_quartic: float = 0.0
    w_solar_headroom: float = 0.0
\end{lstlisting}


        \section{主要構造の読み取り}
        主要クラスは UpperCostConfig。 主要関数は to\_dict, load\_upper\_cost\_config, pick, quad\_penalty, upper\_stage\_cost, upper\_terminal\_cost, active\_upper\_cost\_terms。

        \subsection{主要クラス・関数}
            \begin{longtable}{p{0.07\linewidth} p{0.16\linewidth} p{0.25\linewidth} p{0.40\linewidth}}
            \toprule
            行 & 種別 & 名前 & 補足 \\
            \midrule
            8 & class & \path|UpperCostConfig| &  \\
44 & function & \path|to_dict| & args: self \\
48 & function & \path|_cfg_value| & args: cfg, key, default \\
54 & function & \path|load_upper_cost_config| & args: mpc\_cfg \\
65 & function & \path|pick| & args: name, default \\
110 & function & \path|quad_penalty| & args: x, cap \\
118 & function & \path|upper_stage_cost| & args: cfg \\
263 & function & \path|upper_terminal_cost| & args: cfg \\
283 & function & \path|active_upper_cost_terms| & args: cfg, threshold \\
            \bottomrule
            \end{longtable}

        \subsection{ファイルを上から読んだときの定義順}
            \begin{longtable}{p{0.07\linewidth} p{0.84\linewidth}}
            \toprule
            行 & 内容 \\
            \midrule
            8 & クラス UpperCostConfig を定義する。 \\
48 & 関数 \_cfg\_value を定義する。 \\
54 & 関数 load\_upper\_cost\_config を定義する。 \\
110 & 関数 quad\_penalty を定義する。 \\
118 & 関数 upper\_stage\_cost を定義する。 \\
263 & 関数 upper\_terminal\_cost を定義する。 \\
283 & 関数 active\_upper\_cost\_terms を定義する。 \\
            \bottomrule
            \end{longtable}

        \subsection{import 群の詳細}
            \begin{longtable}{p{0.07\linewidth} p{0.33\linewidth} p{0.50\linewidth}}
            \toprule
            行 & import 文 & このファイルで読む理由 \\
            \midrule
            1 & \path|from __future__ import annotations| & 型ヒントの遅延評価を有効にし、自己参照型や forward reference を安全に扱うため。 このファイル内での使用位置は少ないか、間接利用である。 \\
3 & \path|from dataclasses import asdict, dataclass| & dataclasses から asdict, dataclass を読み込み、このファイルの処理を組み立てるため。 このファイル内での主な使用位置は L7, L45。 \\
4 & \path|from typing import Any, Dict, Optional| & typing から Any, Dict, Optional を読み込み、このファイルの処理を組み立てるため。 このファイル内での主な使用位置は L44, L48, L55, L57, L126, L132, L145, L148, ...。 \\
            \bottomrule
            \end{longtable}

        \section{このファイルを読む前に必要な基礎知識}
        次の章は、構文やROS用語を既知と仮定しないための説明である。

        \subsection{プログラム、プロセス、メモリ、オブジェクトを区別する}
ソースファイルはディスク上の文字列であり、それ自体は走っていない。Python実行可能プログラムがソースを読み、OSがその実行を一つのプロセスとして管理する。プロセスには仮想メモリ、開いているファイル、スレッド、終了コードなどが対応する。


\begin{lstlisting}
ソースファイル -> Pythonインタプリタが読み込む -> OS上のプロセス -> Pythonオブジェクトをメモリに生成 -> 関数やcallbackを実行
\end{lstlisting}


「メモリ上に生成する」とは、実行中プロセスが使う記憶領域に、その値の型、属性、参照関係を表す実体を用意することである。変数は実体そのものというより、そのオブジェクトを指す名前として理解するとPythonの代入が読みやすい。


\begin{lstlisting}[language=python]
node = MPCNode()
alias = node
# nodeとaliasは同じオブジェクトを参照する。
\end{lstlisting}


一つのプロセスには複数スレッドを持たせられる。スレッドは同じプロセスのメモリを共有するため、受信callbackと最適化callbackが同じself.zを書き換える場合は実行順と排他を検討する必要がある。

\paragraph{根拠資料}
\begin{itemize}[leftmargin=1.6em]
\item \href{https://docs.python.org/3/tutorial/classes.html}{Python公式チュートリアル: Classes}
\end{itemize}

\subsection{名前、代入、型、参照}
Pythonの代入文は、右辺を先に評価し、その結果のオブジェクトへ左辺の名前を結び付ける。`x = f()`では、まずfを呼び、その戻り値が得られてからxが更新される。


`float(...)`、`int(...)`、`bool(...)`は型名を呼び出して値を変換する。外部CSV、YAML、ROS parameterから来た値は型が期待どおりとは限らないため、このリポジトリでは明示変換が多い。


`None`は値が無いことを示す単一のオブジェクト、`math.nan`は浮動小数点値ではあるが有効な数値ではないことを示す。両者は用途が異なるため、`is None`と`math.isfinite`を使い分ける。

\paragraph{根拠資料}
\begin{itemize}[leftmargin=1.6em]
\item \href{https://docs.python.org/3/tutorial/classes.html}{Python公式チュートリアル: Classes}
\end{itemize}

\subsection{関数、引数、戻り値、スコープ}
`def`は関数本体をその場で実行する文ではない。関数オブジェクトを作り、その名前を現在の名前空間へ登録する。`f`は関数そのもの、`f()`はその関数を呼んだ結果である。


仮引数は関数定義側の受け取り名、実引数は呼出側から渡す値である。位置引数、キーワード引数、既定値付き引数、`*args`、`**kwargs`は渡し方が異なる。


\begin{lstlisting}[language=python]
def clip_speed(value, minimum=0.0, maximum=110.0):
    return min(maximum, max(minimum, value))

v = clip_speed(120.0, maximum=90.0)
\end{lstlisting}


関数内で作った通常の名前はローカルスコープに属する。メソッドが`self.z`を更新すればオブジェクトの状態が残るが、単に`z`を更新しただけならその呼出しのローカル値である。

\paragraph{根拠資料}
\begin{itemize}[leftmargin=1.6em]
\item \href{https://docs.python.org/3/tutorial/classes.html}{Python公式チュートリアル: Classes}
\end{itemize}

\subsection{module、package、import、console\_scripts}
Pythonファイルはmoduleとして読み込める。複数moduleをディレクトリにまとめたものがpackageである。`from .model import SolarCarModel`の先頭の点は、現在と同じpackage内のmodel moduleを指す相対importである。


import時にはファイルのトップレベル文が上から一度実行される。`def`や`class`は関数・クラスを登録するが、その本体の通常処理は呼び出すまで走らない。


setuptoolsの`console\_scripts`は、端末で使う実行可能名と`package.module:function`を対応付けるインストール時メタデータである。生成される実行用ラッパーはmoduleをimportし、指定関数を呼び、戻り値を終了コードとして扱う小さな入口である。


\begin{lstlisting}
ROS launchのexecutable名 -> install済みconsole script -> mpc_solarcar.mpc_nodeをimport -> main()を呼ぶ
\end{lstlisting}

\paragraph{根拠資料}
\begin{itemize}[leftmargin=1.6em]
\item \href{https://setuptools.pypa.io/en/latest/userguide/entry_point.html}{setuptools公式: Entry Points / Console Scripts}
\item \href{https://docs.python.org/3/tutorial/classes.html}{Python公式チュートリアル: Classes}
\end{itemize}

\subsection{例外、try/finally、with、資源解放}
例外は通常の戻り値とは別経路で異常を呼出元へ伝える。`try/except`は想定した異常を処理し、`finally`は成功・失敗にかかわらず後始末を行う。


`with open(...) as f:`はcontext managerを使い、ブロックを出るとファイルを閉じる。CSVやログの破損を避けるため、開いた資源の所有者と閉じる場所を明確にする。


`except Exception: pass`はノードを止めない利点がある一方、入力異常を隠して原因追跡を難しくする。安全に関係する値では、少なくとも頻度制限付き警告、異常カウンタ、fallback状態のpublishを検討する。



\subsection{class、独自クラス、インスタンス、self、継承}
クラスは、保持するデータと、そのデータを扱う関数を一つの型としてまとめる設計単位である。「独自クラス」とはPythonやROS 2が最初から用意した型ではなく、このプロジェクトが目的に合わせて定義した新しい型を指す。


`class MPCNode(Node)`は、rclpyのNodeを基底クラスとして継承し、Nodeが持つpublisher、subscription、timer、parameterなどの機能にMPC固有機能を追加する。丸括弧内は関数引数ではなく基底クラス指定である。


`MPCNode()`はクラスオブジェクトを呼び出して新しいインスタンスを作る。Pythonは新しい実体を用意し、その実体を第1引数selfとして`\_\_init\_\_`へ渡す。`self`は予約語ではなく慣習的な名前だが、この慣習を守ることでコードの意味が共有される。


\begin{lstlisting}[language=python]
class Counter:
    def __init__(self):
        self.value = 0

    def increment(self):
        self.value += 1

counter = Counter()
counter.increment()
\end{lstlisting}


`super().\_\_init\_\_(...)`は継承元の初期化処理を呼ぶ。MPCNodeではこれを省くとROSノードとして必要な内部実体が作られず、`create\_publisher`などを正しく使えない。

\paragraph{根拠資料}
\begin{itemize}[leftmargin=1.6em]
\item \href{https://docs.python.org/3/tutorial/classes.html}{Python公式チュートリアル: Classes}
\item \href{https://docs.ros.org/en/humble/Tutorials/Beginner-CLI-Tools/Understanding-ROS2-Nodes/Understanding-ROS2-Nodes.html}{ROS 2 Humble公式: Understanding nodes}
\end{itemize}

\subsection{先頭アンダースコア、\_\_init\_\_、名前付けの作法}
`self.\_step\_solar`の先頭1個のアンダースコアは「クラス内部で使う実装詳細」という慣習であり、アクセスを強制禁止する機構ではない。外部コードから呼べるが、公開APIとして依存しない意思表示である。


`\_\_init\_\_`の前後2個のアンダースコアはPythonが定めた特殊メソッド名である。クラス生成時に自動で呼ばれる。任意の名前に前後2個を付けて独自仕様を作ることは避ける。


`\_`で始まるローカル変数は、未使用または外部へ見せない意図を示す場合がある。例えば`\_base\_csv`は戻り値として受け取る必要はあるが、その後の処理では使わないことを読者へ示す。

\paragraph{根拠資料}
\begin{itemize}[leftmargin=1.6em]
\item \href{https://docs.python.org/3/tutorial/classes.html}{Python公式チュートリアル: Classes}
\end{itemize}

\subsection{デコレータと@dataclass}
`@名前`は、直後に定義した関数またはクラスを別の関数へ渡し、その結果で定義名を置き換えるデコレータ構文である。`@Class`という一般構文があるのではなく、`@dataclass`など実際のデコレータ名を書く。


`@dataclass`は型注釈付きフィールドから`\_\_init\_\_`、`\_\_repr\_\_`、比較処理などを自動生成する。物理パラメータや解析結果のように、名前付きデータを一まとまりで運ぶ用途に向く。


\begin{lstlisting}[language=python]
from dataclasses import dataclass

@dataclass
class State:
    soc: float
    temperature_c: float

state = State(soc=0.8, temperature_c=30.0)
\end{lstlisting}


自動生成は検証を自動で保証しない。単位、許容範囲、相互依存制約は`\_\_post\_init\_\_`や別の検証関数で確認する必要がある。

\paragraph{根拠資料}
\begin{itemize}[leftmargin=1.6em]
\item \href{https://docs.python.org/3/library/dataclasses.html}{Python公式ライブラリ: dataclasses}
\item \href{https://docs.python.org/3/tutorial/classes.html}{Python公式チュートリアル: Classes}
\end{itemize}

\subsection{制御、状態、入力、モデル予測制御MPC}
制御対象の内部を表す状態をx、操作入力をu、外乱・予報をwとすると、離散モデルは`x[k+1] = f(x[k], u[k], w[k])`と書ける。ソーラーカーではSoC、電池温度、距離、速度などが状態候補、速度目標や駆動トルクが入力候補になる。


\[
\min_{u_0,\ldots,u_{N-1}} \sum_{k=0}^{N-1}\ell(x_k,u_k,w_k)+V_f(x_N)
\]

MPCは現在状態からNステップ先まで予測し、目的関数と制約を満たす入力系列を求める。ただし実際に適用するのは通常先頭入力だけで、次回は新しい実測状態から再び解く。これがreceding horizonである。


予測モデル、目的関数、制約、ホライズン、solver、初期値のどれかが変わると答えも変わる。「MPCを使う」だけでは仕様は決まらず、これらを単位付きで追う必要がある。

\paragraph{根拠資料}
\begin{itemize}[leftmargin=1.6em]
\item \href{https://docs.scipy.org/doc/scipy/reference/generated/scipy.optimize.minimize.html}{SciPy公式: scipy.optimize.minimize}
\end{itemize}

\subsection{目的関数、制約、L-BFGS-B、SHGO、有限grid証明}
数値最適化器は、利用者が与えた目的関数を複数の候補点で評価し、より小さい値を持つ候補を探す。solverが物理を理解するのではなく、物理と運用価値はcost関数へ書かれる。


L-BFGS-Bは変数ごとの上下限を扱える局所最適化法である。初期値の近くの谷へ収束し得るため、非凸問題では複数seedや大域探索と組み合わせる。successがFalseでも有限な候補が返る場合があるため、採用条件をコード側で決める。


SHGOは定めたsamplingと局所最適化を組み合わせる大域最適化法である。有限Cartesian gridの全列挙は、そのgrid上の最良を証明できるが、連続領域全体の最良を自動的に証明しない。資料ではこの証明範囲を区別する。

\paragraph{根拠資料}
\begin{itemize}[leftmargin=1.6em]
\item \href{https://docs.scipy.org/doc/scipy/reference/generated/scipy.optimize.minimize.html}{SciPy公式: scipy.optimize.minimize}
\item \href{https://docs.scipy.org/doc/scipy/reference/generated/scipy.optimize.shgo.html}{SciPy公式: scipy.optimize.shgo}
\end{itemize}

\subsection{車両力学、電力収支、効率map}
\[
F_{\mathrm{aero}}=\frac{1}{2}\rho C_dA(v+v_{\mathrm{wind}})^2,\quad F_{\mathrm{roll}}\approx mgC_{rr}\cos\theta,\quad F_{\mathrm{grade}}=mg\sin\theta
\]

車輪機械powerは概ね`P\_mech = F\_total * v + P\_inertia`で、駆動時と回生時で異なる効率mapを通してDC側powerへ変換する。mapは速度・torqueなどを軸にした実測または同定tableであり、範囲外の補間・clip規則もモデルの一部である。


\[
P_{\mathrm{pack}}=P_{\mathrm{drive,dc}}-P_{\mathrm{regen,dc}}+P_{\mathrm{aux}}-P_{\mathrm{pv}}
\]

符号規約は必ずコードで確認する。本プロジェクトでは正のpack powerを放電側としてSoCを減らす処理が中心であり、発電と回生はpack負荷を下げる方向に働く。





        \section{関数・クラスを上から順に解説}
        \subsection{L8 クラス UpperCostConfig}
            \begin{itemize}[leftmargin=1.6em]
              \item 行範囲: L8--L45
              \item 所属: module直下
              \item 定義: \path|UpperCostConfig(bases=none)|
              \item docstring: 明示されていない。
              \item このブロックが直接呼ぶ主な関数・メソッド: 特に明示されていない。
              \item 戻り値の要点: 戻り値が無いか、return 文が明示されていない。
              \item 主なローカル代入名: 特になし。
              \item 読み取る主なインスタンス属性: 特になし。
              \item 更新する主なインスタンス属性: 特になし。
              \item 明示的に送出する例外: 明示的raiseはない。
              \item 制御構造の規模: 条件分岐 0、ループ 0、try 0
            \end{itemize}
            \paragraph{この定義を読むためのPython構文}
            \begin{itemize}[leftmargin=1.6em]
            \item class文は新しい型を定義する。丸括弧内は継承する基底クラスである。
\item @で始まる行は、定義したクラスを加工するdecoratorである。
            \end{itemize}
            \paragraph{上から順の処理}
            \begin{enumerate}[leftmargin=1.8em]
            \item objective\_mode に 'weighted' を代入する。
\item w\_wait に 1.0 を代入する。
\item w\_travel\_time に 1.0 を代入する。
\item w\_terminal\_soc\_min に 30.0 を代入する。
\item w\_day\_end\_soc\_min に 100000.0 を代入する。
\item w\_soc\_day\_max に 10000.0 を代入する。
\item w\_soc\_day\_track に 0.0 を代入する。
\item w\_speed\_smooth に 30.0 を代入する。
\item w\_dv\_limit に 2.0 を代入する。
\item w\_speed\_limit に 50.0 を代入する。
\item w\_drive\_window に 100000.0 を代入する。
\item w\_current\_sq に 0.01 を代入する。
            \end{enumerate}
            \paragraph{代表コード断片}
            \begin{lstlisting}[language=Python]
class UpperCostConfig:
    objective_mode: str = "weighted"
    w_wait: float = 1.0
    w_travel_time: float = 1.0
    w_terminal_soc_min: float = 30.0
    w_day_end_soc_min: float = 1.0e5
    w_soc_day_max: float = 1.0e4
    w_soc_day_track: float = 0.0
    w_speed_smooth: float = 30.0
    w_dv_limit: float = 2.0
    w_speed_limit: float = 50.0
    w_drive_window: float = 1.0e5
    w_current_sq: float = 0.01
    w_pack_energy: float = 0.0
    w_joule_loss: float = 0.0
    w_aero_energy: float = 0.0
    w_mech_energy: float = 0.0
    w_speed_quartic: float = 0.0
    w_solar_headroom: float = 0.0
    w_progress_lag: float = 0.0
    w_progress_terminal_lag: float = 0.0
    w_kinetic_pos: float = 0.0
    w_pack_power_slew: float = 0.0
    w_temp: float = 5.0
    w_soc_terminal: float = 0.0
    w_soc_floor_barrier: float = 0.0
    w_uncertainty_reserve: float = 0.0
    speed_quartic_scale_kmh: float = 80.0
    progress_lag_deadband_km: float = 0.0
    soc_solar_headroom_max: float = 0.92
    solar_headroom_power_scale_w: float = 1000.0
    soc_floor_barrier_eps: float = 0.01
    reserve_soc_per_hour: float = 0.0
    reserve_soc_max_extra: float = 0.0
    constraint_penalty: float = 1.0e4
...
\end{lstlisting}

\subsection{L44 関数 UpperCostConfig.to\_dict}
            \begin{itemize}[leftmargin=1.6em]
              \item 行範囲: L44--L45
              \item 所属: \path|UpperCostConfig|
              \item 定義: \path|to_dict(self) -> Dict[str, Any]|
              \item docstring: 明示されていない。
              \item このブロックが直接呼ぶ主な関数・メソッド: asdict
              \item 戻り値の要点: asdict(self)
              \item 主なローカル代入名: 特になし。
              \item 読み取る主なインスタンス属性: 特になし。
              \item 更新する主なインスタンス属性: 特になし。
              \item 明示的に送出する例外: 明示的raiseはない。
              \item 制御構造の規模: 条件分岐 0、ループ 0、try 0
            \end{itemize}
            \paragraph{この定義を読むためのPython構文}
            \begin{itemize}[leftmargin=1.6em]
            \item 第1引数selfは、このメソッドを呼んだインスタンス自身を受け取る慣習名である。
\item `->`以後は戻り値の型注釈であり、実行時の値を自動変換する命令ではない。
            \end{itemize}
            \paragraph{上から順の処理}
            \begin{enumerate}[leftmargin=1.8em]
            \item asdict(self) を返す。
            \end{enumerate}
            \paragraph{代表コード断片}
            \begin{lstlisting}[language=Python]
    def to_dict(self) -> Dict[str, Any]:
        return asdict(self)
\end{lstlisting}

\subsection{L48 関数 \_cfg\_value}
            \begin{itemize}[leftmargin=1.6em]
              \item 行範囲: L48--L51
              \item 所属: module直下
              \item 定義: \path|_cfg_value(cfg: Optional[dict], key: str, default)|
              \item docstring: 明示されていない。
              \item このブロックが直接呼ぶ主な関数・メソッド: isinstance
              \item 戻り値の要点: default / cfg[key]
              \item 主なローカル代入名: 特になし。
              \item 読み取る主なインスタンス属性: 特になし。
              \item 更新する主なインスタンス属性: 特になし。
              \item 明示的に送出する例外: 明示的raiseはない。
              \item 制御構造の規模: 条件分岐 1、ループ 0、try 0
            \end{itemize}
            \paragraph{この定義を読むためのPython構文}
            \begin{itemize}[leftmargin=1.6em]
            \item 特別な構文上の補足は少ない。
            \end{itemize}
            \paragraph{上から順の処理}
            \begin{enumerate}[leftmargin=1.8em]
            \item 条件 isinstance(cfg, dict) and key in cfg を判定し、真なら内部処理を行う。
\item   cfg[key] を返す。
\item default を返す。
            \end{enumerate}
            \paragraph{代表コード断片}
            \begin{lstlisting}[language=Python]
def _cfg_value(cfg: Optional[dict], key: str, default):
    if isinstance(cfg, dict) and key in cfg:
        return cfg[key]
    return default
\end{lstlisting}

\subsection{L54 関数 load\_upper\_cost\_config}
            \begin{itemize}[leftmargin=1.6em]
              \item 行範囲: L54--L107
              \item 所属: module直下
              \item 定義: \path|load_upper_cost_config(mpc_cfg: Optional[dict], *, legacy: Optional[dict] = None, default_drive_window: float = 100000.0) -> UpperCostConfig|
              \item docstring: 明示されていない。
              \item このブロックが直接呼ぶ主な関数・メソッド: UpperCostConfig, \_cfg\_value, float, get, isinstance, lower, pick, str, strip
              \item 戻り値の要点: UpperCostConfig(objective\_mode=str(pick('objective\_mode', 'weighted')).strip().lower(), w\_wait=float(pick('w\_wait', 1.0)), w\_travel\_time=float(pick('w\_travel\_time', 1.0)), w\_terminal\_soc\_min=float(pick('w\_terminal\_soc\_min', 30.0)), w\_day\_end\_soc\_min=float(pick('w\_day\_end\_soc\_min', default\_drive\_window)), w\_soc\_day\_max=float(pick('w\_soc\_day\_max', 10000.0)), w\_soc\_day\_track=float(pick('w\_soc\_day\_track', 0.0)), w\_speed\_smooth=float(pick('w\_speed\_smooth', \_cfg\_value(legacy, 'w\_dv', 30.0))), w\_dv\_limit=float(pick('w\_dv\_limit', 2.0)), w\_speed\_limit=float(pick('w\_speed\_limit', 50.0)), w\_drive\_window=float(pick('w\_drive\_window', default\_drive\_window)), w\_current\_sq=float(pick('w\_current\_sq', \_cfg\_value(legacy, 'w\_current', 0.01))), w\_pack\_energy=float(pick('w\_pack\_energy', 0.0)), w\_joule\_loss=float(pick('w\_joule\_loss', 0.0)), w\_aero\_energy=float(pick('w\_aero\_energy', 0.0)), w\_mech\_energy=float(pick('w\_mech\_energy', 0.0)), w\_speed\_quartic=float(pick('w\_speed\_quartic', 0.0)), w\_solar\_headroom=float(pick('w\_solar\_headroom', 0.0)), w\_progress\_lag=float(pick('w\_progress\_lag', 0.0)), w\_progress\_terminal\_lag=float(pick('w\_progress\_terminal\_lag', 0.0)), w\_kinetic\_pos=float(pick('w\_kinetic\_pos', 0.0)), w\_pack\_power\_slew=float(pick('w\_pack\_power\_slew', 0.0)), w\_temp=float(pick('w\_temp', \_cfg\_value(legacy, 'w\_T', 5.0))), w\_soc\_terminal=float(pick('w\_soc\_terminal', 0.0)), w\_soc\_floor\_barrier=float(pick('w\_soc\_floor\_barrier', 0.0)), w\_uncertainty\_reserve=float(pick('w\_uncertainty\_reserve', 0.0)), speed\_quartic\_scale\_kmh=float(pick('speed\_quartic\_scale\_kmh', 80.0)), progress\_lag\_deadband\_km=float(pick('progress\_lag\_deadband\_km', 0.0)), soc\_solar\_headroom\_max=float(pick('soc\_solar\_headroom\_max', 0.92)), solar\_headroom\_power\_scale\_w=float(pick('solar\_headroom\_power\_scale\_w', 1000.0)), soc\_floor\_barrier\_eps=float(pick('soc\_floor\_barrier\_eps', 0.01)), reserve\_soc\_per\_hour=float(pick('reserve\_soc\_per\_hour', 0.0)), reserve\_soc\_max\_extra=float(pick('reserve\_soc\_max\_extra', 0.0)), constraint\_penalty=float(pick('constraint\_penalty', 10000.0))) / default / nested[name] / legacy[name]
              \item 主なローカル代入名: legacy, nested
              \item 読み取る主なインスタンス属性: 特になし。
              \item 更新する主なインスタンス属性: 特になし。
              \item 明示的に送出する例外: 明示的raiseはない。
              \item 制御構造の規模: 条件分岐 3、ループ 0、try 0
            \end{itemize}
            \paragraph{この定義を読むためのPython構文}
            \begin{itemize}[leftmargin=1.6em]
            \item `->`以後は戻り値の型注釈であり、実行時の値を自動変換する命令ではない。
            \end{itemize}
            \paragraph{上から順の処理}
            \begin{enumerate}[leftmargin=1.8em]
            \item legacy に legacy or \{\} の結果を代入する。
\item nested に \{\} の結果を代入する。
\item 条件 isinstance(mpc\_cfg, dict) を判定し、真なら内部処理を行う。
\item   nested に mpc\_cfg.get('upper\_cost', \{\}) or \{\} の結果を代入する。
\item 関数 pick を定義する。
\item UpperCostConfig(objective\_mode=str(pick('objective\_mode', 'weighted')).strip().lower(), w\_wait=float(pick('w\_wait', 1.0)), w\_travel\_time=float(pick('w\_travel\_time', 1.0)), w\_terminal\_soc\_min=float(pick('w\_terminal\_soc\_min', 30.0)), w\_day\_end\_soc\_min=float(pick('w\_day\_end\_soc\_min', default\_drive\_window)), w\_soc\_day\_max=float(pick('w\_soc\_day\_max', 10000.0)), w\_soc\_day\_track=float(pick('w\_soc\_day\_track', 0.0)), w\_speed\_smooth=float(pick('w\_speed\_smooth', \_cfg\_value(legacy, 'w\_dv', 30.0))), w\_dv\_limit=float(pick('w\_dv\_limit', 2.0)), w\_speed\_limit=float(pick('w\_speed\_limit', 50.0)), w\_drive\_window=float(pick('w\_drive\_window', default\_drive\_window)), w\_current\_sq=float(pick('w\_current\_sq', \_cfg\_value(legacy, 'w\_current', 0.01))), w\_pack\_energy=float(pick('w\_pack\_energy', 0.0)), w\_joule\_loss=float(pick('w\_joule\_loss', 0.0)), w\_aero\_energy=float(pick('w\_aero\_energy', 0.0)), w\_mech\_energy=float(pick('w\_mech\_energy', 0.0)), w\_speed\_quartic=float(pick('w\_speed\_quartic', 0.0)), w\_solar\_headroom=float(pick('w\_solar\_headroom', 0.0)), w\_progress\_lag=float(pick('w\_progress\_lag', 0.0)), w\_progress\_terminal\_lag=float(pick('w\_progress\_terminal\_lag', 0.0)), w\_kinetic\_pos=float(pick('w\_kinetic\_pos', 0.0)), w\_pack\_power\_slew=float(pick('w\_pack\_power\_slew', 0.0)), w\_temp=float(pick('w\_temp', \_cfg\_value(legacy, 'w\_T', 5.0))), w\_soc\_terminal=float(pick('w\_soc\_terminal', 0.0)), w\_soc\_floor\_barrier=float(pick('w\_soc\_floor\_barrier', 0.0)), w\_uncertainty\_reserve=float(pick('w\_uncertainty\_reserve', 0.0)), speed\_quartic\_scale\_kmh=float(pick('speed\_quartic\_scale\_kmh', 80.0)), progress\_lag\_deadband\_km=float(pick('progress\_lag\_deadband\_km', 0.0)), soc\_solar\_headroom\_max=float(pick('soc\_solar\_headroom\_max', 0.92)), solar\_headroom\_power\_scale\_w=float(pick('solar\_headroom\_power\_scale\_w', 1000.0)), soc\_floor\_barrier\_eps=float(pick('soc\_floor\_barrier\_eps', 0.01)), reserve\_soc\_per\_hour=float(pick('reserve\_soc\_per\_hour', 0.0)), reserve\_soc\_max\_extra=float(pick('reserve\_soc\_max\_extra', 0.0)), constraint\_penalty=float(pick('constraint\_penalty', 10000.0))) を返す。
            \end{enumerate}
            \paragraph{代表コード断片}
            \begin{lstlisting}[language=Python]
def load_upper_cost_config(
    mpc_cfg: Optional[dict],
    *,
    legacy: Optional[dict] = None,
    default_drive_window: float = 1.0e5,
) -> UpperCostConfig:
    legacy = legacy or {}
    nested = {}
    if isinstance(mpc_cfg, dict):
        nested = mpc_cfg.get("upper_cost", {}) or {}

    def pick(name: str, default):
        if name in nested:
            return nested[name]
        if name in legacy:
            return legacy[name]
        return default

    return UpperCostConfig(
        objective_mode=str(pick("objective_mode", "weighted")).strip().lower(),
        w_wait=float(pick("w_wait", 1.0)),
        w_travel_time=float(pick("w_travel_time", 1.0)),
        w_terminal_soc_min=float(pick("w_terminal_soc_min", 30.0)),
        w_day_end_soc_min=float(pick("w_day_end_soc_min", default_drive_window)),
        w_soc_day_max=float(pick("w_soc_day_max", 1.0e4)),
        w_soc_day_track=float(pick("w_soc_day_track", 0.0)),
        w_speed_smooth=float(pick("w_speed_smooth", _cfg_value(legacy, "w_dv", 30.0))),
        w_dv_limit=float(pick("w_dv_limit", 2.0)),
        w_speed_limit=float(pick("w_speed_limit", 50.0)),
        w_drive_window=float(pick("w_drive_window", default_drive_window)),
        w_current_sq=float(pick("w_current_sq", _cfg_value(legacy, "w_current", 0.01))),
        w_pack_energy=float(pick("w_pack_energy", 0.0)),
        w_joule_loss=float(pick("w_joule_loss", 0.0)),
        w_aero_energy=float(pick("w_aero_energy", 0.0)),
        w_mech_energy=float(pick("w_mech_energy", 0.0)),
...
\end{lstlisting}

\subsection{L65 関数 load\_upper\_cost\_config.pick}
            \begin{itemize}[leftmargin=1.6em]
              \item 行範囲: L65--L70
              \item 所属: \path|load_upper_cost_config|
              \item 定義: \path|pick(name: str, default)|
              \item docstring: 明示されていない。
              \item このブロックが直接呼ぶ主な関数・メソッド: 特に明示されていない。
              \item 戻り値の要点: default / nested[name] / legacy[name]
              \item 主なローカル代入名: 特になし。
              \item 読み取る主なインスタンス属性: 特になし。
              \item 更新する主なインスタンス属性: 特になし。
              \item 明示的に送出する例外: 明示的raiseはない。
              \item 制御構造の規模: 条件分岐 2、ループ 0、try 0
            \end{itemize}
            \paragraph{この定義を読むためのPython構文}
            \begin{itemize}[leftmargin=1.6em]
            \item 特別な構文上の補足は少ない。
            \end{itemize}
            \paragraph{上から順の処理}
            \begin{enumerate}[leftmargin=1.8em]
            \item 条件 name in nested を判定し、真なら内部処理を行う。
\item   nested[name] を返す。
\item 条件 name in legacy を判定し、真なら内部処理を行う。
\item   legacy[name] を返す。
\item default を返す。
            \end{enumerate}
            \paragraph{代表コード断片}
            \begin{lstlisting}[language=Python]
    def pick(name: str, default):
        if name in nested:
            return nested[name]
        if name in legacy:
            return legacy[name]
        return default
\end{lstlisting}

\subsection{L110 関数 quad\_penalty}
            \begin{itemize}[leftmargin=1.6em]
              \item 行範囲: L110--L115
              \item 所属: module直下
              \item 定義: \path|quad_penalty(x: float, cap: float = 1000.0) -> float|
              \item docstring: 明示されていない。
              \item このブロックが直接呼ぶ主な関数・メソッド: 特に明示されていない。
              \item 戻り値の要点: x * x / 0.0
              \item 主なローカル代入名: x
              \item 読み取る主なインスタンス属性: 特になし。
              \item 更新する主なインスタンス属性: 特になし。
              \item 明示的に送出する例外: 明示的raiseはない。
              \item 制御構造の規模: 条件分岐 2、ループ 0、try 0
            \end{itemize}
            \paragraph{この定義を読むためのPython構文}
            \begin{itemize}[leftmargin=1.6em]
            \item `->`以後は戻り値の型注釈であり、実行時の値を自動変換する命令ではない。
            \end{itemize}
            \paragraph{上から順の処理}
            \begin{enumerate}[leftmargin=1.8em]
            \item 条件 x <= 0.0 を判定し、真なら内部処理を行う。
\item   0.0 を返す。
\item 条件 x > cap を判定し、真なら内部処理を行う。
\item   x に cap の結果を代入する。
\item x * x を返す。
            \end{enumerate}
            \paragraph{代表コード断片}
            \begin{lstlisting}[language=Python]
def quad_penalty(x: float, cap: float = 1.0e3) -> float:
    if x <= 0.0:
        return 0.0
    if x > cap:
        x = cap
    return x * x
\end{lstlisting}

\subsection{L118 関数 upper\_stage\_cost}
            \begin{itemize}[leftmargin=1.6em]
              \item 行範囲: L118--L260
              \item 所属: module直下
              \item 定義: \path|upper_stage_cost(cfg: UpperCostConfig, *, dt_wait: float, dt_travel: float, v_kmh: float, v_prev_kmh: float, vmax_local_kmh: float, drive_limits: Optional[tuple], dv_limit_kmhps: float, I_a: float, V_v: float, P_pv_w: float, P_pack_w: float, P_pack_prev_w: Optional[float], P_mech_wheel_w: float, losses_int_w: float, losses_line_w: float, F_aero_n: float, kinetic_step_wh: float, z_next: float, Tb_next_c: float, term_soc_min: float, soc_min: float, soc_max: float, temp_min_c: float, temp_max_c: float, day_end_soc_min: Optional[float], day_end_crossing: bool, soc_day_end_max: float, soc_day_track_target: Optional[float], soc_day_track_tol: float, I_max: float, I_chg_min: float, V_min: float, V_max: float, time_ahead_h: float, progress_lag_km: float = 0.0) -> float|
              \item docstring: 明示されていない。
              \item このブロックが直接呼ぶ主な関数・メソッド: abs, float, max, min, quad\_penalty
              \item 戻り値の要点: J / J
              \item 主なローカル代入名: J, c, d\_pack\_kw, dv, e\_aero\_wh, e\_loss\_wh, e\_mech\_wh, e\_pack\_wh, lag\_err, pv\_kw, reserve\_extra, reserve\_soc, soc\_gap, solar\_scale, speed\_scale, vmax\_kmh, vmin\_kmh
              \item 読み取る主なインスタンス属性: 特になし。
              \item 更新する主なインスタンス属性: 特になし。
              \item 明示的に送出する例外: 明示的raiseはない。
              \item 制御構造の規模: 条件分岐 18、ループ 0、try 0
            \end{itemize}
            \paragraph{この定義を読むためのPython構文}
            \begin{itemize}[leftmargin=1.6em]
            \item `->`以後は戻り値の型注釈であり、実行時の値を自動変換する命令ではない。
            \end{itemize}
            \paragraph{上から順の処理}
            \begin{enumerate}[leftmargin=1.8em]
            \item J に 0.0 の結果を代入する。
\item 条件 dt\_wait > 0.0 を判定し、真なら内部処理を行う。
\item   J を Add で更新する。
\item J を Add で更新する。
\item 条件 cfg.objective\_mode in \{'fastest', 'fastest\_feasible', 'minimum\_time'\} を判定し、真なら内部処理を行う。
\item   c に cfg.constraint\_penalty の結果を代入する。
\item   dv に (v\_kmh - v\_prev\_kmh) / max(dt\_travel, 0.001) の結果を代入する。
\item   条件 dv\_limit\_kmhps > 0.0 を判定し、真なら内部処理を行う。
\item     J を Add で更新する。
\item   条件 drive\_limits is not None を判定し、真なら内部処理を行う。
\item     (vmin\_kmh, vmax\_kmh) に drive\_limits の結果を代入する。
\item     J を Add で更新する。
\item     J を Add で更新する。
\item   J を Add で更新する。
\item   条件 day\_end\_crossing and day\_end\_soc\_min is not None を判定し、真なら内部処理を行う。
\item     J を Add で更新する。
\item     条件 soc\_day\_end\_max > 0.0 を判定し、真なら内部処理を行う。
\item       J を Add で更新する。
\item   J を Add で更新する。
\item   J を Add で更新する。
\item   J を Add で更新する。
\item   J を Add で更新する。
\item   J を Add で更新する。
\item   J を Add で更新する。
\item   J を Add で更新する。
\item   J を Add で更新する。
\item   J を返す。
\item J を Add で更新する。
\item 条件 soc\_day\_track\_target is not None を判定し、真なら内部処理を行う。
\item   J を Add で更新する。
\item   J を Add で更新する。
\item dv に (v\_kmh - v\_prev\_kmh) / max(dt\_travel, 0.001) の結果を代入する。
\item J を Add で更新する。
\item 条件 dv\_limit\_kmhps > 0.0 を判定し、真なら内部処理を行う。
\item   J を Add で更新する。
\item 条件 drive\_limits is not None を判定し、真なら内部処理を行う。
\item   (vmin\_kmh, vmax\_kmh) に drive\_limits の結果を代入する。
\item   J を Add で更新する。
\item   J を Add で更新する。
\item J を Add で更新する。
\item J を Add で更新する。
\item e\_pack\_wh に max(0.0, P\_pack\_w) * dt\_travel / 3600.0 の結果を代入する。
\item e\_loss\_wh に max(0.0, losses\_int\_w + losses\_line\_w) * dt\_travel / 3600.0 の結果を代入する。
\item e\_aero\_wh に max(0.0, F\_aero\_n) * (v\_kmh / 3.6) * dt\_travel / 3600.0 の結果を代入する。
\item e\_mech\_wh に max(0.0, P\_mech\_wheel\_w) * dt\_travel / 3600.0 の結果を代入する。
\item J を Add で更新する。
\item J を Add で更新する。
\item J を Add で更新する。
\item J を Add で更新する。
\item J を Add で更新する。
\item 条件 cfg.w\_speed\_quartic > 0.0 を判定し、真なら内部処理を行う。
\item   speed\_scale に max(1.0, cfg.speed\_quartic\_scale\_kmh) の結果を代入する。
\item   J を Add で更新する。
\item 条件 cfg.w\_pack\_power\_slew > 0.0 and P\_pack\_prev\_w is not None を判定し、真なら内部処理を行う。
\item   d\_pack\_kw に (float(P\_pack\_w) - float(P\_pack\_prev\_w)) / 1000.0 の結果を代入する。
\item   J を Add で更新する。
\item 条件 cfg.w\_progress\_lag > 0.0 を判定し、真なら内部処理を行う。
\item   lag\_err に max(0.0, float(progress\_lag\_km) - float(cfg.progress\_lag\_deadband\_km)) の結果を代入する。
\item   J を Add で更新する。
\item 条件 cfg.w\_solar\_headroom > 0.0 を判定し、真なら内部処理を行う。
\item   solar\_scale に max(1.0, cfg.solar\_headroom\_power\_scale\_w) の結果を代入する。
\item   pv\_kw に max(0.0, P\_pv\_w) / solar\_scale の結果を代入する。
\item   条件 pv\_kw > 0.0 を判定し、真なら内部処理を行う。
\item     J を Add で更新する。
\item 条件 cfg.w\_soc\_floor\_barrier > 0.0 を判定し、真なら内部処理を行う。
\item   soc\_gap に max(float(z\_next) - float(soc\_min), float(cfg.soc\_floor\_barrier\_eps)) の結果を代入する。
\item   J を Add で更新する。
\item 条件 cfg.w\_uncertainty\_reserve > 0.0 and cfg.reserve\_soc\_per\_hour > 0.0 を判定し、真なら内部処理を行う。
\item   reserve\_extra に min(max(0.0, float(cfg.reserve\_soc\_max\_extra)), max(0.0, float(time\_ahead\_h)) * max(0.0, float(cfg.reserve\_soc\_per\_hour))) の結果を代入する。
\item   reserve\_soc に min(float(soc\_max), float(soc\_min) + reserve\_extra) の結果を代入する。
\item   J を Add で更新する。
\item 条件 day\_end\_crossing and day\_end\_soc\_min is not None を判定し、真なら内部処理を行う。
\item   J を Add で更新する。
\item   条件 soc\_day\_end\_max > 0.0 を判定し、真なら内部処理を行う。
\item     J を Add で更新する。
\item c に cfg.constraint\_penalty の結果を代入する。
\item J を Add で更新する。
\item J を Add で更新する。
\item J を Add で更新する。
\item J を Add で更新する。
            \end{enumerate}
            \paragraph{代表コード断片}
            \begin{lstlisting}[language=Python]
def upper_stage_cost(
    cfg: UpperCostConfig,
    *,
    dt_wait: float,
    dt_travel: float,
    v_kmh: float,
    v_prev_kmh: float,
    vmax_local_kmh: float,
    drive_limits: Optional[tuple],
    dv_limit_kmhps: float,
    I_a: float,
    V_v: float,
    P_pv_w: float,
    P_pack_w: float,
    P_pack_prev_w: Optional[float],
    P_mech_wheel_w: float,
    losses_int_w: float,
    losses_line_w: float,
    F_aero_n: float,
    kinetic_step_wh: float,
    z_next: float,
    Tb_next_c: float,
    term_soc_min: float,
    soc_min: float,
    soc_max: float,
    temp_min_c: float,
    temp_max_c: float,
    day_end_soc_min: Optional[float],
    day_end_crossing: bool,
    soc_day_end_max: float,
    soc_day_track_target: Optional[float],
    soc_day_track_tol: float,
    I_max: float,
    I_chg_min: float,
    V_min: float,
...
\end{lstlisting}

\subsection{L263 関数 upper\_terminal\_cost}
            \begin{itemize}[leftmargin=1.6em]
              \item 行範囲: L263--L280
              \item 所属: module直下
              \item 定義: \path|upper_terminal_cost(cfg: UpperCostConfig, *, z_terminal: float, term_soc_min: float, soc_finish_target: float, soc_finish_tol: float = 0.0, progress_terminal_lag_km: float = 0.0) -> float|
              \item docstring: 明示されていない。
              \item このブロックが直接呼ぶ主な関数・メソッド: float, max, quad\_penalty
              \item 戻り値の要点: J
              \item 主なローカル代入名: J, lag\_err
              \item 読み取る主なインスタンス属性: 特になし。
              \item 更新する主なインスタンス属性: 特になし。
              \item 明示的に送出する例外: 明示的raiseはない。
              \item 制御構造の規模: 条件分岐 2、ループ 0、try 0
            \end{itemize}
            \paragraph{この定義を読むためのPython構文}
            \begin{itemize}[leftmargin=1.6em]
            \item `->`以後は戻り値の型注釈であり、実行時の値を自動変換する命令ではない。
            \end{itemize}
            \paragraph{上から順の処理}
            \begin{enumerate}[leftmargin=1.8em]
            \item J に cfg.constraint\_penalty * quad\_penalty(term\_soc\_min - z\_terminal) の結果を代入する。
\item 条件 soc\_finish\_target > 0.0 を判定し、真なら内部処理を行う。
\item   J を Add で更新する。
\item 条件 cfg.w\_progress\_terminal\_lag > 0.0 を判定し、真なら内部処理を行う。
\item   lag\_err に max(0.0, float(progress\_terminal\_lag\_km) - float(cfg.progress\_lag\_deadband\_km)) の結果を代入する。
\item   J を Add で更新する。
\item J を返す。
            \end{enumerate}
            \paragraph{代表コード断片}
            \begin{lstlisting}[language=Python]
def upper_terminal_cost(
    cfg: UpperCostConfig,
    *,
    z_terminal: float,
    term_soc_min: float,
    soc_finish_target: float,
    soc_finish_tol: float = 0.0,
    progress_terminal_lag_km: float = 0.0,
) -> float:
    J = cfg.constraint_penalty * quad_penalty(term_soc_min - z_terminal)
    if soc_finish_target > 0.0:
        J += cfg.w_soc_terminal * quad_penalty(
            z_terminal - (soc_finish_target + max(0.0, soc_finish_tol))
        )
    if cfg.w_progress_terminal_lag > 0.0:
        lag_err = max(0.0, float(progress_terminal_lag_km) - float(cfg.progress_lag_deadband_km))
        J += cfg.w_progress_terminal_lag * quad_penalty(lag_err)
    return J
\end{lstlisting}

\subsection{L283 関数 active\_upper\_cost\_terms}
            \begin{itemize}[leftmargin=1.6em]
              \item 行範囲: L283--L288
              \item 所属: module直下
              \item 定義: \path|active_upper_cost_terms(cfg: UpperCostConfig, threshold: float = 1e-06) -> Dict[str, float]|
              \item docstring: 明示されていない。
              \item このブロックが直接呼ぶ主な関数・メソッド: abs, float, items, startswith, to\_dict
              \item 戻り値の要点: out
              \item 主なローカル代入名: key, out, value
              \item 読み取る主なインスタンス属性: 特になし。
              \item 更新する主なインスタンス属性: 特になし。
              \item 明示的に送出する例外: 明示的raiseはない。
              \item 制御構造の規模: 条件分岐 1、ループ 1、try 0
            \end{itemize}
            \paragraph{この定義を読むためのPython構文}
            \begin{itemize}[leftmargin=1.6em]
            \item `->`以後は戻り値の型注釈であり、実行時の値を自動変換する命令ではない。
            \end{itemize}
            \paragraph{上から順の処理}
            \begin{enumerate}[leftmargin=1.8em]
            \item out に \{\} の結果を代入する。
\item cfg.to\_dict().items() を順に走査し、各要素を (key, value) に入れて処理する。
\item   条件 key.startswith('w\_') and abs(float(value)) > threshold を判定し、真なら内部処理を行う。
\item     out[key] に float(value) の結果を代入する。
\item out を返す。
            \end{enumerate}
            \paragraph{代表コード断片}
            \begin{lstlisting}[language=Python]
def active_upper_cost_terms(cfg: UpperCostConfig, threshold: float = 1.0e-6) -> Dict[str, float]:
    out = {}
    for key, value in cfg.to_dict().items():
        if key.startswith("w_") and abs(float(value)) > threshold:
            out[key] = float(value)
    return out
\end{lstlisting}











        \section{処理の流れ}
        \begin{enumerate}[leftmargin=1.7em]
        \item ソース中の主要関数を通じて処理を進める。
        \end{enumerate}

        \section{このファイルの位置づけ}
        この資料群では、\path|mpc_solarcar/upper_cost.py| を
        \textbf{planner helper}の中核として扱う。
        したがって、ここで扱う処理は単独で完結せず、
        前段の profile / launch / telemetry と後段の planner / logger / report へ繋がる。

        \end{document}
